% Academic Research Paper — ACM SIGCONF Style Format
\documentclass[10pt,twocolumn]{article}

\usepackage[margin=0.75in]{geometry}
\usepackage{amsmath,amssymb}
\usepackage{booktabs}
\usepackage{cite}
\usepackage[dvipsnames]{xcolor}
\usepackage{hyperref}

\hypersetup{colorlinks=true,linkcolor=Maroon,citecolor=Maroon,urlcolor=Maroon}

\begin{document}

% ZONE:TITLE_AUTHORS:START
\title{\textbf{Verified Document Composition via Modular Agent Specialization}}
\author{
  Morgan Vance \\
  ACM Member \\
  School of Informatics, University of Edinburgh \\
  \texttt{m.vance@ed.ac.uk}
  \and
  Casey Jordan \\
  Department of Computer Science \\
  ETH Z{\"u}rich \\
  \texttt{casey.jordan@inf.ethz.ch}
}
\date{}
\maketitle
% ZONE:TITLE_AUTHORS:END

% ZONE:ABSTRACT:START
\begin{abstract}
Interactive document synthesis tools often produce uncompilable source code when processing multi-turn human revisions. In this paper, we introduce a verification-driven architecture for ACM-compliant manuscript production. We decompose the generation space into isolated AST zones guarded by type checkers and constraint solvers. Our evaluation demonstrates a 99.1\% compilation success rate across 500 benchmark ACM papers with negligible overhead.
\end{abstract}
% ZONE:ABSTRACT:END

% ZONE:KEYWORDS:START
\paragraph{\textbf{CCS Concepts}}
$\bullet$ \textbf{Software and its engineering} $\rightarrow$ Software architectures; Software verification and validation; $\bullet$ \textbf{Computing methodologies} $\rightarrow$ Artificial intelligence.

\paragraph{\textbf{Keywords}}
Document synthesis, multi-agent systems, automated repair, ACM format.
% ZONE:KEYWORDS:END

% ZONE:INTRODUCTION:START
\section{Introduction}
Modern research software engineering mandates publication-grade documentation and manuscripts adhering to strict conference formatting criteria. Generating such documents via generative AI frequently violates macro and formatting contracts, requiring costly manual intervention.

To solve this, we contribute:
\begin{itemize}
  \item A formal zone-based partitioning strategy for ACM publications.
  \item An automated runtime repair pipeline mitigating macro conflicts.
  \item An open-source benchmark suite for evaluating document synthesis.
\end{itemize}
% ZONE:INTRODUCTION:END

% ZONE:METHODOLOGY:START
\section{Methodology}
Our architecture coordinates specialized agents via an immutable AST registry.

\subsection{Constraint Formulation}
Let $\mathcal{C}$ denote the set of structural invariants required by the target publication venue. For each zone $z \in \mathcal{Z}$, the agent must satisfy:
\begin{equation}
\forall c \in \mathcal{C}, \quad \text{Verify}(z, c) = \text{True}
\end{equation}
If any constraint is violated, the repair engine intervenes before compilation.
% ZONE:METHODOLOGY:END

% ZONE:EXPERIMENTS:START
\section{Evaluation}
We evaluated the verification engine across three benchmark suites.

\begin{table}[ht]
\centering
\caption{Benchmark results for ACM manuscript generation.}
\label{tab:acm_results}
\begin{tabular}{lccc}
\toprule
\textbf{Configuration} & \textbf{Valid AST (\%)} & \textbf{Pass Rate (\%)} & \textbf{Time (s)} \\
\midrule
Unconstrained & 62.4 & 58.1 & 12.3 \\
Rule-Based    & 84.1 & 81.7 & 7.4  \\
Our Approach  & \textbf{99.8} & \textbf{99.1} & \textbf{2.8}  \\
\bottomrule
\end{tabular}
\end{table}
% ZONE:EXPERIMENTS:END

% ZONE:CONCLUSION:START
\section{Conclusion}
We have presented an automated verification framework for ACM-compliant manuscript authoring. Our results confirm that domain-specific agent isolation drastically improves compilation reliability.
% ZONE:CONCLUSION:END

% ZONE:REFERENCES:START
\begin{thebibliography}{99}
\bibitem{knuth1974}
D.~E.~Knuth, ``Structured programming with go to statements,'' \emph{ACM Computing Surveys}, vol.~6, no.~4, pp.~261--301, 1974.

\bibitem{dijkstra1968}
E.~W.~Dijkstra, ``Letters to the editor: Go to statement considered harmful,'' \emph{Communications of the ACM}, vol.~11, no.~3, pp.~147--148, 1968.
\end{thebibliography}
% ZONE:REFERENCES:END

\end{document}
